\documentclass[11pt]{amsart}

\usepackage[T1]{fontenc}
\usepackage{lmodern}
\usepackage{microtype}
\usepackage{amsmath,amssymb}
\usepackage{booktabs}
\usepackage[hidelinks]{hyperref}

\hypersetup{
  pdftitle={Reduce-Max Exceeds Backlog 2 on Ten Bamboos}
}

\newtheorem{theorem}{Theorem}
\newtheorem{corollary}[theorem]{Corollary}
\theoremstyle{remark}
\newtheorem{remark}[theorem]{Remark}

\newcommand{\RM}{\textnormal{\textsc{Reduce-Max}}}
\newcommand{\state}[2]{\bigl[#1\mid #2\bigr]}

\title{Reduce-Max Exceeds Backlog~$2$ on Ten Bamboos}
\date{}
\subjclass[2020]{68W40}
\keywords{bamboo garden trimming, greedy algorithm, backlog, counterexample}

\begin{document}

\begin{abstract}
After normalizing total daily growth to~$1$, D'Emidio, Di Stefano,
and Navarra conjectured that the greedy bamboo-trimming rule \RM,
which cuts a tallest bamboo in each round, has backlog at most~$2$.
Jasi\'nska, Kuszmaul, and Lee disproved this conjecture with instances
on $2000$ and $2702$ bamboos, and report that all of their
counterexamples have between $1200$ and $3000$ bamboos and that
similar steep distributions much smaller or larger yielded none.
We give an exact ten-bamboo
witness for which every tie-breaking rule has backlog
\[
\frac{1878}{937}=2+\frac{4}{937}>2.
\]
The maximum is attained in round~$60$, and from the post-round-$137$
state onward the trajectory modulo permutation of six equal-rate
bamboos is seven-periodic. The verification uses only integer
arithmetic. As a consequence, positive-rate counterexamples exist for
every number of bamboos $n\ge 10$.
\end{abstract}

\maketitle

\section{The counterexample}

A bamboo instance has positive growth rates $r_1,\dots,r_n$ with
$\sum_i r_i=1$. Initially all heights are zero. In each round, every
bamboo grows by its rate and one bamboo of maximum height is then cut
to height zero. The backlog is the supremum of the maximum pre-cut
height over all rounds. This greedy rule is called \RM.

D'Emidio, Di Stefano, and Navarra conjectured that \RM\ is
$2$-approximating with respect to total daily growth; under the
normalization above, this is the backlog bound~$2$
\cite[Conjecture~1]{DDN19}. Here $2$-approximating is used in the
sense of \cite{DDN19}, where the property observed in their
experiments is recorded as $M(E)<2H$ with $H$ the sum of the growth
rates; it is that inequality which fails below. Nothing here bears on
the ratio of $M(E)$ to the optimal backlog of the individual instance.
Jasi\'nska, Kuszmaul, and Lee disproved
the conjecture, giving a $2000$-bamboo example with backlog $2.0004$
and a $2702$-bamboo example yielding the stronger lower bound
$2.076$~\cite[Section~4]{JKL26}. That disproof is theirs and no
priority for it is claimed here: the contribution below is an exact
witness with ten bamboos, and it does not improve the numerical lower
bound $2.076$. What it does settle is a question of size left open
in~\cite{JKL26}, where it is reported that ``all of our
counter-examples have between $1200$ and $3000$ bamboos'', that
``Examining similar steep distributions on much smaller or larger
distributions did not yield examples with backlog greater than $2$'',
and that this is ``likely because for smaller number of bamboos, the
slower bamboos will not be numerous enough''~\cite[Section~4.1]{JKL26}.
Theorem~\ref{thm:main} shows that ten bamboos suffice, two orders of
magnitude below that range, so no such lower limit on the number of
bamboos holds.

\begin{theorem}\label{thm:main}
For the rate vector
\[
 r=\frac{1}{937}
 \bigl(313,261,131,118,19,19,19,19,19,19\bigr),
\]
every tie-breaking of \RM\ has backlog exactly $1878/937$.
\end{theorem}

\begin{proof}
The numerator weights
\[
w=(313,261,131,118,19,19,19,19,19,19)
\]
sum to $937$. Scale all heights by $937$. Let $H(t)$ be the integer
post-cut state after round $t$, with $H(0)=0$, and set
\[
\widehat H_i(t)=H_i(t-1)+w_i.
\]
If bamboo $j_t$ is cut in round $t$, then
$H_{j_t}(t)=0$ and $H_i(t)=\widehat H_i(t)$ for $i\ne j_t$.
Let $C(H)$ retain the first four coordinates and sort the last six in
nondecreasing order.

A symbol $k\in\{1,2,3,4\}$ means that bamboo $k$ is cut. A symbol $S$
means that a bamboo of maximum pre-cut height among coordinates
$5,\dots,10$ is cut. Since these six bamboos have the same rate, the
resulting canonical state is independent of the choice among tied
slow bamboos. For the indicated cut $j_t$, define the cross-rate gap
\[
\delta_t=
\widehat H_{j_t}(t)-
\max_{\,w_i\ne w_{j_t}}\widehat H_i(t).
\]
The following table is an exact certificate. The columns $m$, $M$,
and $\Delta$ are, respectively, the minimum and maximum selected
pre-cut heights and the minimum value of $\delta_t$ on the stated
interval.

\begin{table}[ht]
\centering
\caption{Integer certificate for the greedy trajectory.}
\label{tab:certificate}
\begin{tabular}{@{}cclrrr@{}}
\toprule
Rounds & Length & Cut word & $m$ & $M$ & $\Delta$ \\
\midrule
$1$--$6$     & $6$  & $121321$             & $313$ & $939$  & $2$  \\
$7$--$42$    & $36$ & $(421321)^6$         & $708$ & $939$  & $4$  \\
$43$--$60$   & $18$ & $SS214321SS21SS4321$ & $817$ & $1878$ & $1$  \\
$61$--$72$   & $12$ & $(214321)^2$         & $522$ & $1252$ & $3$  \\
$73$--$100$  & $28$ & $(S214321)^4$        & $570$ & $1252$ & $23$ \\
$101$--$108$ & $8$  & $SS214321$           & $874$ & $1252$ & $27$ \\
$109$--$137$ & $29$ & $(S214321)^4S$       & $684$ & $1252$ & $40$ \\
$138$--$144$ & $7$  & $214321S$            & $783$ & $1252$ & $40$ \\
\bottomrule
\end{tabular}
\end{table}

For reference, the canonical post-cut states at the interval
endpoints are
\begin{align*}
C(H(0))   &=\state{0,0,0,0}{0,0,0,0,0,0},\\
C(H(6))   &=\state{0,261,262,708}{114,114,114,114,114,114},\\
C(H(42))  &=\state{0,261,262,590}{798,798,798,798,798,798},\\
C(H(60))  &=\state{0,261,262,354}{76,95,152,171,304,323},\\
C(H(72))  &=\state{0,261,262,354}{304,323,380,399,532,551},\\
C(H(100)) &=\state{0,261,262,354}{114,247,380,513,836,855},\\
C(H(108)) &=\state{0,261,262,354}{114,133,266,399,532,665},\\
C(H(137)) &=\state{313,522,393,472}{0,133,266,399,532,665},\\
C(H(144)) &=\state{313,522,393,472}{0,133,266,399,532,665}.
\end{align*}
All entries follow directly from the recurrence above. The positive
values of $\Delta$ show that every indicated cut is globally greedy:
for a digit, its rate class is a singleton; for $S$, the selected
bamboo is maximal within the slow class by definition. The only ties
for the global maximum occur in rounds
\[
43,\ 44,\ 51,\ 52,\ 55,
\]
and each is among equal-rate, equal-height slow bamboos. Hence every
tie-breaking induces the same canonical trajectory.

The decisive comparisons in rounds $56$--$60$ (the multipliers count
rounds since the most recent cut) are
\begin{align*}
56:&\quad 56\cdot19=1064>9\cdot118=1062,\\
57:&\quad 10\cdot118=1180>9\cdot131=1179,\\
58:&\quad 10\cdot131=1310>5\cdot261=1305,\\
59:&\quad 6\cdot261=1566>5\cdot313=1565,\\
60:&\quad 6\cdot313=1878=2\cdot937+4.
\end{align*}
Thus round~$60$ has pre-cut height $1878/937>2$. Table
\ref{tab:certificate} shows that no other round through $144$ has a
larger height. Moreover, $C(H(137))=C(H(144))$, and the intervening
cut word is $214321S$. The canonical greedy evolution therefore
repeats with period~$7$, with cycle peak $1252/937$. Consequently the
global backlog is exactly $1878/937$, independently of tie-breaking.
\end{proof}

\begin{corollary}\label{cor:alln}
For every integer $n\ge10$, there is a positive-rate $n$-bamboo
instance for which every tie-breaking of \RM\ has backlog greater
than $2$.
\end{corollary}

\begin{proof}
Write $n=10+m$ with $m\ge0$. The case $m=0$ is
Theorem~\ref{thm:main}. For $m\ge1$, multiply the ten rates in the
theorem by $999/1000$ and add $m$ bamboos, each of rate
$1/(1000m)$. The rates are positive and sum to~$1$.

Through round~$60$, every selected bamboo in the original execution
has pre-cut height at least $313/937$, by
Table~\ref{tab:certificate}. Inductively, before any new bamboo has
been cut, its height in a round $t\le60$ is at most
\[
\frac{t}{1000m}\le\frac{60}{1000}
<\frac{999}{1000}\frac{313}{937}.
\]
Thus no new bamboo can affect the canonical evolution of the original
ten during these rounds. In round~$60$, the scaled first bamboo
reaches
\[
\frac{999}{1000}\frac{1878}{937}
=2+\frac{2122}{937000}>2.
\]
The conclusion holds for every tie-breaking because the original
ten-bamboo canonical trajectory is tie-independent.
\end{proof}

\begin{remark}
The theorem gives a substantially smaller exact witness than the two
examples displayed in~\cite{JKL26}, but it neither improves their
lower bound $2.076$ nor proves minimality. In particular, no claim is
made about the existence of a counterexample with at most nine
bamboos. The instance also lies outside the exhaustive search that
supported \cite[Conjecture~1]{DDN19}. Those experiments enumerated
integer partitions of the total daily growth $H$, for
$H\in\{5,10,\dots,35\}$
\cite[Sections~4.1 and~4.2]{DDN19}. What excludes the rate vector of
Theorem~\ref{thm:main} is the bound on $H$: in integer form it has
$H=937$, so it was never enumerated there, even though its $n$ is
only~$10$.
\end{remark}

\begin{thebibliography}{9}

\bibitem{DDN19}
M.~D'Emidio, G.~Di Stefano, and A.~Navarra,
\emph{Bamboo garden trimming problem: priority schedulings},
Algorithms \textbf{12} (2019), no.~4, Article~74.
\href{https://doi.org/10.3390/a12040074}
{doi:10.3390/a12040074}.

\bibitem{JKL26}
K.~Jasi\'nska, J.~Kuszmaul, and G.~Lee,
\emph{Strengths and limitations of greedy in cup games},
arXiv:2602.18610v1 [cs.DS], 2026.
\href{https://doi.org/10.48550/arXiv.2602.18610}
{doi:10.48550/arXiv.2602.18610}.

\end{thebibliography}

\end{document}